\begin{theorem}[Algorithmic fixed-DFA prime-slice obstruction]
  \label{thm:effective-fixed-dfa-prime-obstruction}
  Fix \(b\ge2\) and a complete DFA
  \(\mathcal A=(Q,\Sigma_b,\delta,q_0,F)\), encoded by its integer
  transition table. Let \(B\) be its unlabeled transition matrix and
  \(P=B/b\). There is an algorithm which outputs the following
  compressed exact certificate.
  \begin{enumerate}[label=\textup{(\alph*)}]
    \item The transient set \(T\), the recurrent classes
    \(R_1,\dots,R_s\) of \(P\), their periods \(d_j\), and cyclic
    decompositions \(R_{j,0},\dots,R_{j,d_j-1}\).
    \item For every \(j,t\), the stationary probability vector
    \(\pi_{j,t}\) of the primitive return chain
    \(P^{d_j}|_{R_{j,t}}\), with rational entries.
    \item The absorption probabilities \(H_j(q)\), \(q\in T\), obtained
    from
    \[
      (I-P_{TT})H_j=P_{T,R_j}\mathbf 1,
    \]
    and the refined cyclic absorption probabilities
    \[
      A_{j,t}(q)
      :=\sum_{\ell\ge0}
        e_q^{\mathsf T}P_{TT}^{\ell}P_{T,R_j}
        S_{j,t+\ell+1}\mathbf 1 ,
    \]
    where indices are taken modulo \(d_j\) and
    \(S_{j,u}\) is the cyclic projection onto \(R_{j,u}\).
    Thus, conditional on absorption after \(\ell\) transient steps, the
    recurrent entry class is chosen so that at every later time
    \(m\equiv r\pmod {d_j}\) the recurrent chain lies in
    \(R_{j,t+r}\). The displayed series is evaluated exactly as a
    finite family of linear solves after grouping \(\ell\) modulo
    \(d_j\), namely
    \[
      A_{j,t}
      =
      \sum_{u=0}^{d_j-1}
      (I-P_{TT}^{d_j})^{-1}P_{TT}^uP_{T,R_j}
      S_{j,t+u+1}\mathbf 1 .
    \]
    \item For each residue \(r\) modulo
    \(D_A=\operatorname{lcm}_j d_j\), the compressed formula
    \[
      c_r=
      \sum_{j=1}^s\sum_{t=0}^{d_j-1}
      \alpha_{j,t}^{(r)}\,
      \pi_{j,t+r}(F\cap R_{j,t+r}),
      \tag{\(*\)}
      \label{eq:compressed-cr}
    \]
    where
    \[
      \alpha_{j,t}^{(r)}=
      \begin{cases}
        1, & q_0\in R_{j,t},\\
        A_{j,t}(q_0), & q_0\in T,\\
        0, & \text{otherwise.}
      \end{cases}
    \]
    Again all cyclic indices are read modulo \(d_j\). This formula
    decides \(c_r=0\) or \(c_r>0\) by exact rational arithmetic.
    \item Rational constants \(0<\tau_r<\theta_r<1\), \(K_r>0\), and
    integer thresholds \(N_r\) such that
    \[
      \abs{a_m(\mathcal A)/b^m-c_r}\le K_r\theta_r^m
      \qquad(m\equiv r\pmod {D_A},\ m\ge N_r).
    \]
    Here \(\tau_r\) is the certified decay rate for the transient and
    return-error blocks, while \(\theta_r\) is the larger final
    exponential rate after the polynomial convolution factor has been
    absorbed. The certificate records the strict spectral-radius
    inequalities and rational resolvent bounds of
    \Cref{lem:effective-spectral-constants}, including the rational
    numbers \(\tau_r,\theta_r\), rational upper bounds \(K_r\), the
    monotone tail threshold, and the finite list of exact rational
    inequalities checked before the tail. These data are enough to
    verify the inequality mechanically, without reconstructing an
    algebraic Jordan basis.
  \end{enumerate}
  \smallskip
  \noindent\textit{Certificate schema.}
  For each residue class the certificate schema consists of the
  following fields, all represented by rational data except for the
  root-isolation rectangles.
  \begin{enumerate}[label=\textup{(\roman*)}]
    \item Rational matrices \(P,T,Q_{j,t},E_{j,t}=Q_{j,t}-\Pi_{j,t}\),
    the cyclic projections \(S_{j,t}\), and the rational vectors
    \(\pi_{j,t},A_{j,t}\) solving the displayed linear systems.
    The norm convention throughout the schema is the
    \(\ell^\infty\) operator norm, i.e. maximum absolute row sum.
    \item Rational numbers \(c_r,\tau_r,\theta_r,K_r,C_r\), integers
    \(J_r,N_r\), and a tail threshold \(N_r^{\mathrm{tail}}\), with
    \(0<\tau_r<\theta_r<1\).
    \item Root-isolation certificates for every characteristic
    polynomial: rational rectangles, each containing exactly one complex
    root and disjoint from the others, whose squared-modulus intervals
    certify \(\rho(T)<\tau_r\) and
    \(\rho(E_{j,t})<\tau_r^{d_j}\). The verifier checks these by
    Sturm--subdivision or any equivalent exact rational root-counting
    procedure on the rectangles and by rational interval arithmetic for
    the squared moduli.
    \item Resolvent certificates
    \((\gamma_A,\delta_A,U_A,L_A)\) for
    \(A\in\{T\}\cup\{E_{j,t}\}_{j,t}\), with
    \(\rho(A)<\gamma_A<\tau_r\) for \(A=T\) and
    \(\rho(A)<\gamma_A<\tau_r^{d_j}\) for \(A=E_{j,t}\). The verifier
    checks, using the same isolating rectangles and rational interval
    arithmetic, that every rectangle lies strictly inside the circle
    \(|z|=\gamma_A\) and supplies a positive rational distance
    \(\Delta_\xi\) from the rectangle to that circle. Therefore
    \[
      |\det(zI-A)|\ge \prod_{\xi}\Delta_\xi=:\delta_A
      \qquad(|z|=\gamma_A),
    \]
    with roots counted with algebraic multiplicity. The verifier also
    checks the elementary coefficient bound
    \[
      \|\operatorname{adj}(zI-A)\|_\infty\le U_A,
    \]
    and hence that \(\|A^n\|_\infty\le L_A\gamma_A^n\) for all
    \(n\ge0\).
    \item A rational block-form expansion of \(P^m-S_r\) on the residue
    class \(m\equiv r\pmod {D_A}\). For every recurrent block \(R_j\)
    and residue \(a\bmod d_j\) the data include the rational cyclic step
    matrices \(G_{j,u}=P_{R_{j,u},R_{j,u+1}}\), their products
    \(W_{j,t,a}=G_{j,t}\cdots G_{j,t+a-1}\), the rational limit matrix
    \(C_{j,a}^{\infty}\), and the checked identities
    \[
      C_j^{d_jh+a}-C_{j,a}^{\infty}
      =
      \sum_t M_{j,t,a}(h)
    \]
    for \(h\ge1\). Here \(M_{j,t,a}(h)\) is zero except on the
    \(R_{j,t}\)-to-\(R_{j,t+a}\) block, where it equals
    \(W_{j,t,a}E_{j,t+a}^h\). The finitely many \(h=0\) cases are
    checked directly. For every block from the transient part to \(R_j\),
    the data also include the exact identity
    \[
      (P^m-S_r)_{T,R_j}
      =
      \sum_{\ell=0}^{m-1}
        T^\ell P_{T,R_j}
        \bigl(C_j^{m-1-\ell}-C_{j,r-\ell-1}^{\infty}\bigr)
      -
      \sum_{\ell=m}^{\infty}
        T^\ell P_{T,R_j}C_{j,r-\ell-1}^{\infty}.
    \]
    These are finite rational identities: the second infinite sum is
    verified after grouping \(\ell\bmod D_A\) as a finite list of
    geometric resolvents \((I-T^{D_A})^{-1}T^u\). The verifier checks
    each displayed identity by rational matrix multiplication and exact
    inversion over \(\QQ\). For the first finite convolution, the
    verifier further splits \(\ell=u+D_Aq\). If the recurrent factor in
    that residue has length \(m-1-\ell=d_jh+a\), with
    \(0\le a<d_j\), then
    \[
      h=H_{m,u}-L_jq,\qquad
      L_j:=D_A/d_j,\qquad
      H_{m,u}:=(m-1-u-a)/d_j .
    \]
    The integer \(a\) is fixed by the listed residue class, so
    \(H_{m,u}\) is integral and the return-error exponent decreases as
    \(q\) increases. The certificate records only the terms with
    \(H_{m,u}-L_jq\ge1\); the finitely many smaller recurrent exponents
    are moved to the direct finite-check list. The resulting finite
    residue list is
    \[
      \sum_{q=0}^{Q(m,u)}
        X_{u}T^{D_Aq}Y_{u}E_{j,t_u}^{H_{m,u}-L_jq}Z_{u},
    \]
    with all matrices \(X_u,Y_u,Z_u\) rational. This is the normal form
    for every convolution term.
    \item The exact convolution bounds attached to the preceding
    expansion. For every direct term \(X A^nY\) the certificate lists
    the rational product
    \(\|X\|_\infty L_A\|Y\|_\infty\eta^n\). For every convolution term
    it lists the rational coefficient
    \[
      C_u=\|X_u\|_\infty L_T\|Y_u\|_\infty
        L_{E_{j,t_u}}\|Z_u\|_\infty\tau_r^{-1-u-a}
    \]
    for the corresponding residue \(u,a\); if \(a=0\), this means the
    same displayed formula with the factor \(\tau_r^{-1-u}\). Indeed
    \(T^{D_Aq}\) contributes at most \(\tau_r^{D_Aq}\), while
    \(E_{j,t_u}^{H_{m,u}-L_jq}\) contributes at most
    \(\tau_r^{d_jH_{m,u}-D_Aq}\); the \(q\)-dependence cancels, and the
    identity \(m-1-u=d_jH_{m,u}+a\) gives
    \[
      \tau_r^{D_Aq}\tau_r^{d_jH_{m,u}-D_Aq}
      =\tau_r^{m-1-u-a},
    \]
    which is bounded by the displayed \(C_u\tau_r^m\) coefficient for
    each \(q\). The number of terms is at most \(m\),
    so the residue summand is bounded by \(C_um\tau_r^m\). For every
    tail it lists the
    rational coefficient in
    \(C\,\tau_r^m/(1-\tau_r^{D_A})\). Summing these finitely many
    coefficients gives the stated coefficient bound
    proving
    \[
      \|P^m-S_r\|_\infty\le C_rm^{J_r}\tau_r^m
      \qquad(m\equiv r\pmod {D_A}).
    \]
    \item The rational tail inequality
    \[
      C_rm^{J_r}\tau_r^m\le K_r\theta_r^m
      \qquad(m\ge N_r^{\mathrm{tail}})
    \]
    and exact rational matrix-power checks for every
    \(m<N_r^{\mathrm{tail}}\) in the residue class that is not covered
    by the tail. The displayed exponential estimate is accepted with
    \(N_r\) no smaller than both the finite-check cutoff and
    \(N_r^{\mathrm{tail}}\).
  \end{enumerate}
  \smallskip
  \noindent\textit{Certificate-verifier theorem.}
  Given the fields just listed, there is a finite verifier which accepts
  if and only if the following rational assertions all check:
  the stochastic block decomposition is a permutation of \(P\); each
  \(C_{j,a}^{\infty}\), \(\pi_{j,t}\), and \(A_{j,t}\) satisfies its
  displayed linear equations; the root-isolation rectangles prove the
  strict spectral-radius inequalities; the resolvent certificates prove
  the listed bounds for every decaying block; the finite block normal
  form for \(P^m-S_r\) is exactly the one displayed in the preceding
  two items; the summed direct, convolution, and tail coefficients give
  \(\|P^m-S_r\|_\infty\le C_rm^{J_r}\tau_r^m\) on the residue class;
  the finitely many pre-tail powers satisfy the same advertised
  estimate after enlargement of \(K_r\); and the rational monotone-tail
  inequality gives \(C_rm^{J_r}\tau_r^m\le K_r\theta_r^m\). All steps
  use rational matrix arithmetic, exact inverses over \(\QQ\), finite
  root counts in rational rectangles, and rational interval arithmetic.
  No unlisted block formula or limiting argument is part of the
  verification step.
  With this output, the following holds for every
  \(m\ge M_r\) satisfying \(m\equiv r\pmod {D_A}\).
  If \(c_r=0\), then
  \[
    a_m(\mathcal A)\le K_r\theta_r^m b^m
    \quad\text{and}\quad
    \abs{(L(\mathcal A)\cap\Sigma_b^m)\triangle\mathcal P_m^{(b)}}
    \ge \frac{A_b}{2}\,\frac{b^m}{m}.
  \]
  If \(c_r>0\), then
  \[
    a_m(\mathcal A)\ge \frac{c_r}{2}b^m,
    \qquad
    \mathrm{Prec}_m\le \frac{2B_b}{c_r}\,\frac1m,
  \]
  whenever \(\mathrm{Prec}_m\) is defined, i.e. on lengths where the
  denominator is nonzero, and
  \[
    \abs{(L(\mathcal A)\cap\Sigma_b^m)\triangle\mathcal P_m^{(b)}}
    \ge \frac{c_r}{4} b^m .
  \]
  Here \(A_b,B_b,M_b\) may be the explicit constants supplied by
  \Cref{lem:base-prime-count}. Consequently
  one computes constants \(c_{\mathcal A,b}>0\) and \(M_{\mathcal A,b}\)
  such that
  \[
    \abs{(L(\mathcal A)\cap\Sigma_b^m)\triangle\mathcal P_m^{(b)}}
    \ge c_{\mathcal A,b}\,\frac{b^m}{m}
    \qquad(m\ge M_{\mathcal A,b}).
  \]
  The compressed output is the tuple
  \[
    \mathcal C=(D_A,\{R_j,d_j,(R_{j,t})_t,\pi_{j,t},A_{j,t}\}_{j,t},
      \{c_r,\tau_r,\theta_r,K_r,C_r,J_r,N_r,M_r\}_{r\bmod D_A})
  \]
  when all residues are enumerated. If \(D_A\) is large, the same data
  may instead be represented by the class list together with an exact
  residue-evaluation procedure for
  \(r\mapsto c_r,\tau_r,\theta_r,K_r,C_r,J_r,N_r,M_r\) and the
  corresponding proof object from part~\textup{(e)};
  any final all-large constant obtained by taking a minimum over
  residues is then represented by this finite minimization procedure.
  Printing the full residue table is charged to output size \(D_A\).
  No polynomial bound in the input size alone is asserted for the
  spectral certificate, for enumerating all residues, or for the final
  minimization.
\end{theorem}

\begin{proof}
  Compute the matrix \(B=\sum_{d\in\Sigma_b}T_d\), the stochastic matrix
  \(P=B/b\), the transient set, the recurrent classes, their periods,
  their cyclic decompositions, and the common period \(D_A\) by standard
  graph algorithms. The vectors \(\pi_{j,t}\) are the unique solutions
  of
  \[
    \pi_{j,t}P^{d_j}|_{R_{j,t}}=\pi_{j,t},
    \qquad
    \pi_{j,t}\mathbf 1=1.
  \]
  Because \(P\) has rational entries, these vectors are rational and
  are obtained by rational Gaussian elimination. The absorption
  probabilities \(H_j\) solve the displayed systems in the statement;
  the matrix \(I-P_{TT}\) is invertible because \(T\) is transient.
  For the cyclic refinements, group the defining series for
  \(A_{j,t}\) by \(\ell\equiv u\pmod{d_j}\). Each group is a
  rational expression in
  \[
    (I-P_{TT}^{d_j})^{-1}P_{TT}^uP_{T,R_j}S_{j,t+u+1}\mathbf 1,
  \]
  hence is again computed by exact rational linear algebra. Formula
  \eqref{eq:compressed-cr} is precisely the canonical block-power limit
  of \(u^{\mathsf T}P^m v\) along \(m\equiv r\pmod {D_A}\). Indeed,
  suppose the chain starts in a transient state \(q\), remains
  transient for exactly \(\ell\) steps, and then enters \(R_j\) at time
  \(\ell+1\) in cyclic class \(R_{j,u}\). Because \(P\) maps
  \(R_{j,a}\) into \(R_{j,a+1}\), at a later time
  \(m\equiv r\pmod {d_j}\) its recurrent cyclic class is
  \[
    u+m-(\ell+1)\equiv u+r-\ell-1 \pmod {d_j}.
  \]
  To contribute to the limiting distribution
  \(\pi_{j,t+r}\) on \(R_{j,t+r}\), the entry class must therefore
  satisfy \(u+r-\ell-1\equiv t+r\), equivalently
  \(u\equiv t+\ell+1\pmod {d_j}\). This is exactly the projection
  \(S_{j,t+\ell+1}\) in the definition of \(A_{j,t}\). Conditioning on
  the absorption time and then applying the primitive return-chain limit
  on \(R_{j,t+r}\) proves \eqref{eq:compressed-cr}.

  Applying \Cref{lem:effective-spectral-constants} to \(P\) and
  \(D_A\) gives \(\tau_r,\theta_r,K_r,N_r\) and
  \[
    \abs{a_m(\mathcal A)/b^m-c_r}\le K_r\theta_r^m
    \qquad(m\equiv r\bmod {D_A},\ m\ge N_r).
  \]
  The associated certificate consists of rational
  \(\tau_r,\theta_r,K_r,C_r\), the integer \(J_r\), the
  spectral-radius proof objects for the transient and return-error
  blocks, the rational resolvent bounds, the rational inequality
  \(C_rm^{J_r}\tau_r^m\le K_r\theta_r^m\) on the tail, and the finite
  rational power checks preceding the tail. Thus the quantified
  inequality used below is part of the output, not an unverifiable
  analytic side condition.

  Choose \(A_b,B_b>0\) and \(M_b\) from
  \Cref{lem:base-prime-count}, so that
  \[
    A_b\,\frac{b^m}{m}
    \le\abs{\mathcal P_m^{(b)}}
    \le B_b\,\frac{b^m}{m}
    \qquad(m\ge M_b).
  \]
  If \(c_r=0\), take \(M_r\ge \max(M_b,N_r)\) to be any integer in the
  residue class tail satisfying
  \[
    K_r\theta_r^m\le \frac{A_b}{2m}
    \qquad(m\ge M_r,\ m\equiv r\bmod {D_A}).
  \]
  Since \(K_r\) and \(\theta_r\) are rational, this threshold is
  obtained by checking finitely many values before monotonicity of
  \(mK_r\theta_r^m\) begins and then recording the first value after
  which the monotone tail is below \(A_b/2\).
  For \(m\ge M_r\) with \(m\equiv r\bmod {D_A}\),
  \[
    a_m(\mathcal A)\le \frac{A_b}{2}\frac{b^m}{m},
  \]
  so the symmetric difference is at least
  \(A_b b^m/(2m)\).

  If \(c_r>0\), take \(M_r\ge \max(M_b,N_r)\) so that
  \[
    K_r\theta_r^m\le c_r/2,
    \qquad
    B_b/m\le c_r/4
  \]
  for all \(m\ge M_r\) with \(m\equiv r\bmod {D_A}\). This is again a
  rational monotone-tail check.
  Then \(a_m(\mathcal A)\ge(c_r/2)b^m\), whence
  \[
    \mathrm{Prec}_m
    \le\frac{\abs{\mathcal P_m^{(b)}}}{a_m(\mathcal A)}
    \le\frac{2B_b}{c_r}\frac1m,
  \]
  and
  \[
    \abs{(L(\mathcal A)\cap\Sigma_b^m)\triangle\mathcal P_m^{(b)}}
    \ge a_m(\mathcal A)-\abs{\mathcal P_m^{(b)}}
    \ge \frac{c_r}{4}b^m.
  \]
  This proves the positive-density lower bound.

  Taking
  \(M_{\mathcal A,b}:=\max_r M_r\) and the minimum over the finitely many
  positive constants \(A_b/2\) and \(c_r/4\) appearing above gives the
  final all-large-\(m\) constant after replacing \(c_r b^m/4\) by the
  weaker \(c_r b^m/(4m)\) in the positive-density residue classes.
  The output statement follows from strongly connected component and
  period algorithms, exact rational linear algebra, root isolation for
  the finitely many characteristic polynomials, and the certified
  resolvent construction in \Cref{lem:effective-spectral-constants}.
  The certificate is compressed by recurrent class and cyclic class.
  Printing all residues, and taking the final minimum by enumeration,
  are charged to the chosen output representation.
\end{proof}

\begin{example}[Two explicit certificates]
  \label{ex:base-ending-zero-certificate}
  Let \(\mathcal A_0\) be the complete DFA over \(\Sigma_b\) accepting
  exactly the non-empty words ending in \(0\). It has states
  \(q_0,q_1\), initial state \(q_0\), accepting set \(\{q_1\}\), and
  transitions
  \[
    \delta(q_i,0)=q_1,\qquad
    \delta(q_i,d)=q_0\quad(1\le d\le b-1).
  \]
  Its unlabeled transition matrix and stochastic matrix are
  \[
    B=\begin{pmatrix} b-1&1\\ b-1&1\end{pmatrix},
    \qquad
    P=\begin{pmatrix} (b-1)/b&1/b\\ (b-1)/b&1/b\end{pmatrix}.
  \]
  There is one recurrent aperiodic class \(R_1=\{q_0,q_1\}\), no
  transient state, \(d_1=D_A=1\), and
  \[
    \pi_{1,0}=\bigl((b-1)/b,\;1/b\bigr),\qquad
    c_0=\pi_{1,0}(\{q_1\})=1/b.
  \]
  Since \(P^2=P\), one may take \(S_0=P\), \(N_0=1\), any rational
  \(0<\tau_0<\theta_0<1\), and any rational \(K_0>0\) for the advertised
  convergence estimate on all \(m\ge1\). Equivalently,
  \(a_m(\mathcal A_0)=b^{m-1}\) and
  \(a_m(\mathcal A_0)/b^m=1/b\) for every \(m\ge1\).

  Let \(A_b,B_b,M_b\) be the constants from
  \Cref{lem:base-prime-count}. The positive-density branch of
  \Cref{thm:effective-fixed-dfa-prime-obstruction} may therefore use the
  explicit threshold
  \[
    M_0:=\left\lceil\max\{M_b,4bB_b\}\right\rceil .
  \]
  For every \(m\ge M_0\),
  \[
    \mathrm{Prec}_m
    \le
    \frac{\abs{\mathcal P_m^{(b)}}}{a_m(\mathcal A_0)}
    \le
    \frac{bB_b}{m},
  \]
  on lengths where the denominator is nonzero, and
  \[
    \abs{(L(\mathcal A_0)\cap\Sigma_b^m)\triangle\mathcal P_m^{(b)}}
    \ge
    b^{m-1}-B_b\frac{b^m}{m}
    \ge
    \frac{1}{2b}b^m .
  \]
  Thus the certificate has \(D_A=1\), \(c_0=1/b\), precision constant
  \(bB_b\), symmetric-difference constant \(1/(2b)\) on the stronger
  \(b^m\) scale, and the final theorem also permits the weaker constant
  \(1/(2b)\) on the \(b^m/m\) scale.

  The empty DFA gives the complementary sparse branch. Here
  \(a_m=0\), \(c_0=0\), and any rational \(K_0>0\) is admissible. Hence, for all
  \(m\ge M_b\),
  \[
    \abs{(\varnothing\cap\Sigma_b^m)\triangle\mathcal P_m^{(b)}}
    =
    \abs{\mathcal P_m^{(b)}}
    \ge
    A_b\frac{b^m}{m}.
  \]
\end{example}

We first establish the quantitative obstruction that drives both
corollaries below.

\begin{proposition}[Quantitative recall/precision obstruction]
  \label{prop:dfa-quantitative-obstruction}
  Set
  \[
    a_m(\mathcal A):=\abs{L(\mathcal A)\cap\Sigma_b^m},
  \]
  and define
  \[
    \mathrm{Rec}_m
    :=\frac{\abs{L(\mathcal A)\cap\mathcal P_m^{(b)}}}{\abs{\mathcal P_m^{(b)}}},
    \qquad
    \mathrm{Prec}_m
    :=\frac{\abs{L(\mathcal A)\cap\mathcal P_m^{(b)}}}{a_m(\mathcal A)}
  \]
  whenever \(a_m(\mathcal A)>0\). Suppose that there exist \(\delta>0\)
  and infinitely many integers \(m\) such that
  \(\mathrm{Rec}_m\ge \delta\) and \(\mathrm{Prec}_m\ge \delta\). Then
  on those lengths,
  \[
    \delta\,\abs{\mathcal P_m^{(b)}}
    \le \abs{L(\mathcal A)\cap\mathcal P_m^{(b)}}
    \le a_m(\mathcal A)
    \le \delta^{-1}\abs{L(\mathcal A)\cap\mathcal P_m^{(b)}}
    \le \delta^{-1}\abs{\mathcal P_m^{(b)}}.
  \]
  In particular, \(a_m(\mathcal A)\asymp b^m/m\) along infinitely many
  lengths, contradicting \Cref{thm:dfa-density-dichotomy}.
\end{proposition}

\begin{proof}
  The lower bound follows from
  \[
    \abs{L(\mathcal A)\cap\mathcal P_m^{(b)}}
    =\mathrm{Rec}_m\,\abs{\mathcal P_m^{(b)}}
    \ge \delta\,\abs{\mathcal P_m^{(b)}},
  \]
  and the upper bound follows from
  \[
    a_m(\mathcal A)
    =\frac{\abs{L(\mathcal A)\cap\mathcal P_m^{(b)}}}{\mathrm{Prec}_m}
    \le \delta^{-1}\abs{L(\mathcal A)\cap\mathcal P_m^{(b)}}
    \le \delta^{-1}\abs{\mathcal P_m^{(b)}}.
  \]
  By \Cref{lem:base-prime-count}, the displayed chain gives
  \(a_m(\mathcal A)\asymp b^m/m\) on an infinite set \(M\) of lengths.

  Let \(D_A\), \(c_0,\dots,c_{D_A-1}\), and \(\theta\) be supplied by
  \Cref{thm:dfa-density-dichotomy}. By the pigeonhole principle, some
  residue class \(r\bmod D_A\) contains infinitely many lengths from
  \(M\). Along this infinite subsequence the preceding bounds give
  positive constants \(K_1,K_2\) such that
  \[
    \frac{K_1}{m}\le \frac{a_m(\mathcal A)}{b^m}
    \le \frac{K_2}{m}.
  \]
  If \(c_r=0\), then the density dichotomy gives
  \(a_m(\mathcal A)/b^m=O(\theta^m)\) on \(m\equiv r\pmod {D_A}\), which is
  incompatible with \(K_1/m\) on an infinite subsequence. If \(c_r>0\),
  then \(a_m(\mathcal A)/b^m\ge c_r/2\) for all sufficiently large
  \(m\equiv r\pmod {D_A}\), which is incompatible with the upper bound
  \(K_2/m\). Both residue-class alternatives contradict the assumed
  infinite set \(M\).
\end{proof}

\begin{corollary}[Symmetric-difference lower bound for base-\(b\) primes]
  \label{cor:dfa-prime-symmetric-diff}
  Let \(L(\mathcal A)\subseteq\Sigma_b^\ast\) be recognised by a fixed
  DFA \(\mathcal A\). Then there exists \(c_{\mathcal A,b}>0\) such that
  \[
    \abs{(L(\mathcal A)\cap\Sigma_b^m)\triangle \mathcal P_m^{(b)}}
    \ge c_{\mathcal A,b}\,\frac{b^m}{m}
  \]
  for all sufficiently large \(m\). If \(a_m(\mathcal A)/b^m\) lies in
  the positive-density alternative of \Cref{thm:dfa-density-dichotomy},
  then along an infinite arithmetic progression of lengths one has the
  stronger bound
  \[
    \abs{(L(\mathcal A)\cap\Sigma_b^m)\triangle \mathcal P_m^{(b)}}\gg b^m.
  \]
\end{corollary}

\begin{proof}
  Let \(D_A\ge 1\), \(c_0,\dots,c_{D_A-1}\ge 0\), and \(\theta\in(0,1)\) be
  as in \Cref{thm:dfa-density-dichotomy}. Replacing \(\theta\) by
  \(\theta^{D_A}\) if necessary, we may write for each residue class
  \(r\bmod D_A\)
  \[
    \frac{a_{D_Ak+r}(\mathcal A)}{b^{D_Ak+r}}=c_r+O(\theta^k)
    \qquad (k\to\infty).
  \]
  Fix constants \(A_b,B_b>0\), supplied by
  \Cref{lem:base-prime-count}, such that
  \[
    A_b\,\frac{b^m}{m}\le \abs{\mathcal P_m^{(b)}}\le B_b\,\frac{b^m}{m}
  \]
  for all sufficiently large \(m\).

  We now treat each residue class separately.

  If \(c_r=0\), then
  \[
    a_{D_Ak+r}(\mathcal A)=O(\theta^k b^{D_Ak+r})
    =o\!\left(\frac{b^{D_Ak+r}}{D_Ak+r}\right).
  \]
  Hence there exists \(k_r\) such that for all \(k\ge k_r\),
  \[
    a_{D_Ak+r}(\mathcal A)
    \le \frac{A_b}{2}\,\frac{b^{D_Ak+r}}{D_Ak+r}.
  \]
  For those \(k\),
  \[
    \abs{(L(\mathcal A)\cap\Sigma_b^{D_Ak+r})\triangle \mathcal P_{D_Ak+r}^{(b)}}
    \ge \abs{\mathcal P_{D_Ak+r}^{(b)}}-a_{D_Ak+r}(\mathcal A)
    \ge \frac{A_b}{2}\,\frac{b^{D_Ak+r}}{D_Ak+r}.
  \]

  If \(c_r>0\), then there exists \(k_r\) such that for all \(k\ge k_r\),
  \[
    a_{D_Ak+r}(\mathcal A)\ge \frac{c_r}{2}\,b^{D_Ak+r}.
  \]
  Enlarging \(k_r\) if necessary, we may also ensure that
  \(B_b/(D_Ak+r)\le c_r/4\). Then for all such \(k\),
  \[
    \abs{(L(\mathcal A)\cap\Sigma_b^{D_Ak+r})\triangle \mathcal P_{D_Ak+r}^{(b)}}
    \ge a_{D_Ak+r}(\mathcal A)-\abs{\mathcal P_{D_Ak+r}^{(b)}}
    \ge \frac{c_r}{4}\,b^{D_Ak+r}.
  \]

  Since there are only finitely many residue classes, taking the minimum
  of \(A_b/2\) and the finitely many numbers \(c_r/4\) arising in the
  case \(c_r>0\) yields a constant \(c_{\mathcal A,b}>0\) such that
  \[
    \abs{(L(\mathcal A)\cap\Sigma_b^m)\triangle \mathcal P_m^{(b)}}
    \ge c_{\mathcal A,b}\,\frac{b^m}{m}
  \]
  for all sufficiently large \(m\). If \(a_m(\mathcal A)/b^m\) falls in
  the positive-density alternative of \Cref{thm:dfa-density-dichotomy},
  then some residue class has \(c_r>0\), and along that infinite
  arithmetic progression the stronger bound
  \(\abs{(L(\mathcal A)\cap\Sigma_b^m)\triangle \mathcal P_m^{(b)}}\gg b^m\)
  holds.
\end{proof}

\begin{remark}[Sharpness of the lower bounds]\label{rem:prime-diff-sharp}
  The two scales in \Cref{cor:dfa-prime-symmetric-diff} are optimal up
  to the implicit constants. If \(L=\emptyset\), then
  \[
    \abs{(L(\mathcal A)\cap\Sigma_b^m)\triangle\mathcal P_m^{(b)}}
    =\abs{\mathcal P_m^{(b)}}\asymp b^m/m .
  \]
  On the other hand, for the positive-density language
  \(\Sigma_b^\ast0\) the symmetric difference with
  \(\mathcal P_m^{(b)}\) is \(\asymp b^m\) along every residue class,
  so the strengthened \(\Omega(b^m)\) conclusion on arithmetic
  progressions cannot be improved.
\end{remark}

\begin{corollary}[No simultaneous recall and precision for fixed DFA]
  \label{cor:dfa-prime-recall-precision}
  With \(\mathrm{Rec}_m\) and \(\mathrm{Prec}_m\) as in
  \Cref{prop:dfa-quantitative-obstruction}, let \(D_A\) and
  \(c_0,\dots,c_{D_A-1}\) be the automaton period and residue constants supplied by
  \Cref{thm:dfa-density-dichotomy}. For each residue class
  \(r\bmod D_A\) exactly one of the following alternatives holds:
  \begin{enumerate}[label=\textup{(\alph*)}]
    \item If \(c_r=0\), then \(\mathrm{Rec}_m\) decays exponentially
    along \(m\equiv r\pmod {D_A}\).
    \item If \(c_r>0\), then there is a constant
    \(C_r=2B_b/c_r\), where
    \(\abs{\mathcal P_m^{(b)}}\le B_b b^m/m\) for all large \(m\),
    such that \(\mathrm{Prec}_m\le C_r/m\) for all sufficiently large
    \(m\equiv r\pmod {D_A}\) for which \(a_m(\mathcal A)>0\).
  \end{enumerate}
  Consequently every DFA satisfies one of the following statements:
  \begin{enumerate}[label=\textup{(\roman*)}]
    \item \(\mathrm{Rec}_m\to 0\) exponentially fast.
    \item There exists at least one residue class
    \(r\bmod D_A\) with \(c_r>0\) such that
    \(\mathrm{Prec}_m\le C_{\mathcal A,b}/m\) for all sufficiently large
    \(m\equiv r\pmod {D_A}\) for which precision is defined.
  \end{enumerate}
  Precision is left undefined on lengths with \(a_m(\mathcal A)=0\);
  no assertion is made about other positive-density or sparse residue
  classes beyond the per-residue alternatives above.
  In particular, no fixed DFA has both recall and precision bounded
  away from zero for infinitely many lengths.
\end{corollary}

\begin{proof}
  Apply \Cref{thm:dfa-density-dichotomy}. Thus there exist an integer
  \(D_A\ge 1\), constants \(c_0,\dots,c_{D_A-1}\ge 0\), and
  \(\theta\in(0,1)\) such that
  \[
    \frac{a_m(\mathcal A)}{b^m}=c_{m\bmod D_A}+O(\theta^m).
  \]
  Fix the upper-bound constant \(B_b>0\) from
  \Cref{lem:base-prime-count}, so that
  \(\abs{\mathcal P_m^{(b)}}\le B_b b^m/m\) for all sufficiently large
  \(m\).
  We treat each residue class \(r\bmod D_A\) separately by considering the
  two cases \(c_r=0\) and \(c_r>0\).

  \textit{Case \(c_r=0\).}
  Along the subsequence \(m\equiv r\pmod {D_A}\), we have
  \(a_m(\mathcal A)=O(\theta^m b^m)\). Since
  \Cref{lem:base-prime-count} gives
  \(\abs{\mathcal P_m^{(b)}}\asymp b^m/m\), so
  \[
    \mathrm{Rec}_m
    \le \frac{a_m(\mathcal A)}{\abs{\mathcal P_m^{(b)}}}
    \ll \theta^m\cdot m.
  \]
  Because \(m\theta^m=O(\theta_r^m)\) for some \(\theta_r\in(0,1)\),
  recall decays exponentially along this residue class.

  \textit{Case \(c_r>0\).}
  Along \(m\equiv r\pmod {D_A}\) with \(m\) large enough, we have
  \(a_m(\mathcal A)\ge \tfrac{1}{2}c_r b^m\). Using the bound for
  \(\abs{\mathcal P_m^{(b)}}\),
  \[
    \mathrm{Prec}_m
    \le \frac{\abs{\mathcal P_m^{(b)}}}{a_m(\mathcal A)}
    \le \frac{2B_b}{c_r}\cdot\frac{1}{m}.
  \]
  Hence precision decays at the rate \(O(1/m)\) along this
  arithmetic progression, so we may take
  \(C_{\mathcal A,b}=2B_b/c_r\) for that subsequence.

  Every residue class \(r\bmod D_A\) falls into one of these two cases.
  If all residue classes have \(c_r=0\), then
  \(\mathrm{Rec}_m\to 0\) exponentially for every \(m\), giving
  alternative~\textup{(i)}. Otherwise, at least one residue class has
  \(c_r>0\), and \(\mathrm{Prec}_m\le C_{\mathcal A,b}/m\) on defined
  lengths along that infinite arithmetic progression, giving
  alternative~\textup{(ii)}.

  If both recall and precision were bounded below by some
  \(\delta>0\) for infinitely many lengths, then
  \Cref{prop:dfa-quantitative-obstruction} would force
  \(a_m(\mathcal A)\asymp b^m/m\) along an infinite set of lengths. But the
  per-residue analysis above shows that on each residue class either
  \(a_m(\mathcal A)\) is exponentially sparse or \(a_m(\mathcal A)\gg b^m\), and
  neither is compatible with \(a_m(\mathcal A)\asymp b^m/m\).
\end{proof}

